\documentclass[11pt]{article}

\usepackage[margin=1in]{geometry}
\usepackage{amsmath,amssymb,amsthm,mathtools}
\usepackage{booktabs}
\usepackage{enumitem}
\usepackage{algorithm}
\usepackage{algpseudocode}

\newcommand{\studentname}{Xilang Chen}
\newcommand{\studentid}{524031910668}

\DeclarePairedDelimiter{\floor}{\lfloor}{\rfloor}
\DeclarePairedDelimiter{\ceil}{\lceil}{\rceil}
\newcommand{\Mod}[1]{\ (\mathrm{mod}\ #1)}

\newtheorem*{fact}{Fact}
\theoremstyle{definition}
\newtheorem*{remark}{Remark}

\setlength{\parindent}{0pt}
\setlength{\parskip}{0.6em}

\title{SE3308 Algorithm Design I --- Assignment 1\\[0.3em]
  \large Exercises 0.1, 0.2, 1.14, 1.20, 1.31, 1.35}
\author{\studentname\ifx\studentid\empty\else\quad\studentid\fi}
\date{}

\begin{document}
\maketitle

Throughout, $\log$ denotes $\log_2$ and $\ln$ the natural logarithm.

%======================================================================
\section*{Exercise 0.1}

We use the following standard facts.

\begin{fact}[Limit rule]
Let $f,g>0$ and suppose $L=\lim_{n\to\infty} f(n)/g(n)$ exists in $[0,\infty]$.
If $0<L<\infty$ then $f=\Theta(g)$; if $L=0$ then $f=O(g)$ but $f\neq\Omega(g)$;
if $L=\infty$ then $f=\Omega(g)$ but $f\neq O(g)$.
\end{fact}

\begin{fact}[Logs vs.\ polynomials]
For any constants $a,b>0$, $\lim_{n\to\infty}(\log n)^b/n^a=0$.
(Substituting $n=2^t$, this is $t^b/2^{at}\to 0$: an exponential beats any polynomial.)
\end{fact}

\begin{center}
\renewcommand{\arraystretch}{1.25}
\begin{tabular}{clll}
\toprule
 & $f(n)$ & $g(n)$ & Answer \\
\midrule
(a) & $n-100$ & $n-200$ & $f=\Theta(g)$ \\
(b) & $n^{1/2}$ & $n^{2/3}$ & $f=O(g)$ \\
(c) & $100n+\log n$ & $n+(\log n)^2$ & $f=\Theta(g)$ \\
(d) & $n\log n$ & $10n\log 10n$ & $f=\Theta(g)$ \\
(e) & $\log 2n$ & $\log 3n$ & $f=\Theta(g)$ \\
(f) & $10\log n$ & $\log(n^2)$ & $f=\Theta(g)$ \\
(g) & $n^{1.01}$ & $n\log^2 n$ & $f=\Omega(g)$ \\
(h) & $n^2/\log n$ & $n(\log n)^2$ & $f=\Omega(g)$ \\
(i) & $n^{0.1}$ & $(\log n)^{10}$ & $f=\Omega(g)$ \\
(j) & $(\log n)^{\log n}$ & $n/\log n$ & $f=\Omega(g)$ \\
(k) & $\sqrt n$ & $(\log n)^3$ & $f=\Omega(g)$ \\
(l) & $n^{1/2}$ & $5^{\log_2 n}$ & $f=O(g)$ \\
(m) & $n2^n$ & $3^n$ & $f=O(g)$ \\
(n) & $2^n$ & $2^{n+1}$ & $f=\Theta(g)$ \\
(o) & $n!$ & $2^n$ & $f=\Omega(g)$ \\
(p) & $(\log n)^{\log n}$ & $2^{(\log_2 n)^2}$ & $f=O(g)$ \\
(q) & $\sum_{i=1}^n i^k$ & $n^{k+1}$ & $f=\Theta(g)$ \\
\bottomrule
\end{tabular}
\end{center}

Wherever the answer is only $O$ or only $\Omega$, the other relation fails (the ratio tends to $0$ or $\infty$).

\textbf{Justifications.}
\begin{enumerate}[label=(\alph*), leftmargin=2.2em, itemsep=0.25em]
\item $\dfrac{n-100}{n-200}\to 1$.
\item $\dfrac{n^{1/2}}{n^{2/3}}=n^{-1/6}\to 0$.
\item $\dfrac{100n+\log n}{n+(\log n)^2}=\dfrac{100+(\log n)/n}{1+(\log n)^2/n}\to 100$.
\item $\dfrac{n\log n}{10n\log(10n)}=\dfrac{\log n}{10(\log n+\log 10)}\to\dfrac1{10}$.
\item $\dfrac{\log 2n}{\log 3n}=\dfrac{1+\log n}{\log 3+\log n}\to 1$.
\item $\log(n^2)=2\log n$, so $f/g=5$ exactly.
\item $\dfrac{n^{1.01}}{n\log^2 n}=\dfrac{n^{0.01}}{(\log n)^2}\to\infty$.
\item $\dfrac{n^2/\log n}{n(\log n)^2}=\dfrac{n}{(\log n)^3}\to\infty$.
\item $\dfrac{n^{0.1}}{(\log n)^{10}}\to\infty$.
\item $(\log n)^{\log n}=2^{\log n\cdot\log\log n}=n^{\log\log n}$. For $n\ge 16$, $\log\log n\ge 2$, so
  $\dfrac{f}{g}=n^{\log\log n-1}\log n\ge n\to\infty$.
\item $\dfrac{\sqrt n}{(\log n)^3}\to\infty$.
\item $5^{\log_2 n}=2^{\log_2 5\cdot\log_2 n}=n^{\log_2 5}$ with $\log_2 5\approx 2.32$, so
  $f/g=n^{1/2-\log_2 5}\to 0$.
\item $\dfrac{n2^n}{3^n}=n\left(\tfrac23\right)^n\to 0$.
\item $2^{n+1}=2\cdot 2^n$, so $f/g=\tfrac12$.
\item $\dfrac{n!}{2^n}=\prod_{i=1}^n\dfrac i2\ge\dfrac12\left(\dfrac32\right)^{n-2}\to\infty$ for $n\ge2$
  (the first factor is $\tfrac12$, the second is $1$, and every later factor is at least $\tfrac32$).
\item As in (j), $f=2^{\log n\cdot\log\log n}$, while $g=2^{(\log n)^2}$. Hence
  $f/g=2^{-\log n\,(\log n-\log\log n)}\to 0$, since $\log n-\log\log n\to\infty$.
\item Here $k\ge 0$ is a fixed constant. Upper bound: each of the $n$ terms is at most $n^k$, so
  $\sum_{i=1}^n i^k\le n^{k+1}$. Lower bound: the terms with $i\ge\ceil{n/2}$ number at least $n/2$
  and each is at least $(n/2)^k$, so $\sum_{i=1}^n i^k\ge (n/2)^{k+1}=2^{-(k+1)}n^{k+1}$.
  Thus $f=\Theta(n^{k+1})$.
\end{enumerate}

\begin{remark}
Changing the base of a logarithm multiplies it by a constant, so none of the answers depend on the base
(in (j) and (p) the logarithm is in the exponent, but the same comparisons go through with $\ln$).
\end{remark}

%======================================================================
\section*{Exercise 0.2}

Let $g(n)=\sum_{i=0}^{n}c^i$ with $c>0$. All terms are positive.

\begin{enumerate}[label=(\alph*), leftmargin=2.2em]
\item $c<1$. Since $g(n)$ contains the term $1$ and all terms are positive, $g(n)\ge 1$. Also
  \[
    g(n)\le\sum_{i=0}^{\infty}c^i=\frac{1}{1-c},
  \]
  a constant independent of $n$. Hence $1\le g(n)\le\frac1{1-c}$ for all $n$, i.e.\ $g(n)=\Theta(1)$.

\item $c=1$. Every term equals $1$, so $g(n)=n+1$. For $n\ge1$, $n\le n+1\le 2n$, hence $g(n)=\Theta(n)$.

\item $c>1$. The last term gives $g(n)\ge c^n$. By the geometric sum formula,
  \[
    g(n)=\frac{c^{n+1}-1}{c-1}<\frac{c^{n+1}}{c-1}=\frac{c}{c-1}\cdot c^n .
  \]
  Since $\frac{c}{c-1}$ is a positive constant, $c^n\le g(n)\le\frac{c}{c-1}c^n$, i.e.\ $g(n)=\Theta(c^n)$.
\end{enumerate}

%======================================================================
\section*{Exercise 1.14}

\textbf{Idea.} Combine the matrix method of Exercise 0.4 with modular repeated squaring.
Let $X=\begin{pmatrix}0&1\\1&1\end{pmatrix}$. Since
$\begin{pmatrix}F_{k+1}\\F_{k+2}\end{pmatrix}=X\begin{pmatrix}F_k\\F_{k+1}\end{pmatrix}$, induction gives
\[
  \begin{pmatrix}F_n\\F_{n+1}\end{pmatrix}=X^n\begin{pmatrix}F_0\\F_1\end{pmatrix}
  =X^n\begin{pmatrix}0\\1\end{pmatrix},
\]
so $F_n$ is the top-right entry $(X^n)_{12}$. It therefore suffices to compute $X^n \bmod p$
(every entry reduced modulo $p$).

\begin{algorithm}[H]
\caption{Computing $F_n \bmod p$}
\begin{algorithmic}[1]
\Function{MatPow}{$n,p$} \Comment{returns $X^n \bmod p$}
  \If{$n=0$} \State \Return $I_2 \bmod p$ \EndIf
  \State $Z\gets$ \Call{MatPow}{$\floor{n/2},p$}
  \State $Z\gets Z\cdot Z \bmod p$
  \If{$n$ is odd} \State $Z\gets X\cdot Z \bmod p$ \EndIf
  \State \Return $Z$
\EndFunction
\Function{FibMod}{$n,p$}
  \State \Return the top-right entry of \Call{MatPow}{$n,p$}
\EndFunction
\end{algorithmic}
\end{algorithm}

Here ``$A\cdot B\bmod p$'' means the $2\times2$ matrix product (8 multiplications and 4 additions of
entries, Exercise 0.4(a)) followed by reducing each of the 4 entries modulo $p$.

\textbf{Correctness.} \textsc{MatPow} follows the recursion
$X^n=(X^{\floor{n/2}})^2$ for even $n$ and $X^n=X\cdot(X^{\floor{n/2}})^2$ for odd $n$, exactly as in
modular exponentiation. Every matrix entry is computed from others using only $+$ and $\times$, and
reduction modulo $p$ commutes with these operations
($a\equiv a',\,b\equiv b'\Rightarrow a+b\equiv a'+b',\ ab\equiv a'b' \pmod p$).
Hence by induction on $n$, each entry of the returned matrix is congruent modulo $p$ to the corresponding
entry of $X^n$, and lies in $\{0,\dots,p-1\}$. In particular \textsc{FibMod} returns $F_n\bmod p$.

\textbf{Running time.} Let $k=\ceil{\log(p+1)}$ be the number of bits of $p$.
\begin{itemize}[itemsep=0.2em]
\item The recursion depth is $\floor{\log n}+1=O(\log n)$, and each level performs at most two matrix
  multiplications, so there are $O(\log n)$ matrix multiplications in total.
\item All entries are in $[0,p)$, i.e.\ at most $k$ bits. One matrix multiplication does 8 products of
  $k$-bit numbers ($O(k^2)$ each), 4 additions of $2k$-bit numbers ($O(k)$ each), and 4 reductions of
  $O(k)$-bit numbers modulo $p$ ($O(k^2)$ each). So one matrix multiplication costs $O(k^2)$.
\end{itemize}
The total running time is
\[
  O(\log n\cdot\log^2 p)
\]
bit operations, which is polynomial in the input size ($\approx\log n+\log p$ bits).
By contrast, running \texttt{fib2} and reducing modulo $p$ at every step takes $n$ additions,
i.e.\ $O(n\log p)$ time, which is exponential in the length $\log n$ of $n$.

%======================================================================
\section*{Exercise 1.20}

We use the extended Euclidean algorithm: $a$ has an inverse modulo $N$ iff $\gcd(a,N)=1$, and if
$ax+Ny=1$ then $x\bmod N$ is the inverse.

\begin{enumerate}[leftmargin=2.2em, itemsep=0.6em]
\item $20^{-1}\bmod 79$.
  \[
    79=3\cdot20+19,\qquad 20=1\cdot19+1 .
  \]
  Back-substituting: $1=20-19=20-(79-3\cdot20)=4\cdot20-1\cdot79$.
  Hence $20^{-1}\equiv \boxed{4}\pmod{79}$. Check: $20\cdot4=80=79+1$.

\item $3^{-1}\bmod 62$.
  \[
    62=20\cdot3+2,\qquad 3=1\cdot2+1 .
  \]
  Back-substituting: $1=3-2=3-(62-20\cdot3)=21\cdot3-1\cdot62$.
  Hence $3^{-1}\equiv \boxed{21}\pmod{62}$. Check: $3\cdot21=63=62+1$.

\item $21^{-1}\bmod 91$.
  \[
    91=4\cdot21+7,\qquad 21=3\cdot7+0 ,
  \]
  so $\gcd(21,91)=7\neq1$ (indeed $21=3\cdot7$, $91=13\cdot7$). Therefore
  \textbf{21 has no inverse modulo 91}: if $21x\equiv1\pmod{91}$, then $7$ would divide
  $21x-91y=1$, which is impossible.

\item $5^{-1}\bmod 23$.
  \[
    23=4\cdot5+3,\qquad 5=1\cdot3+2,\qquad 3=1\cdot2+1 .
  \]
  Back-substituting:
  \[
    1=3-2=3-(5-3)=2\cdot3-5=2(23-4\cdot5)-5=2\cdot23-9\cdot5 .
  \]
  Hence $5^{-1}\equiv-9\equiv \boxed{14}\pmod{23}$. Check: $5\cdot14=70=3\cdot23+1$.
\end{enumerate}

%======================================================================
\section*{Exercise 1.31}

Let $N$ be an $n$-bit number, so $2^{n-1}\le N<2^n$ and $n-1\le\log N<n$.

\textbf{(a)} The number of bits of $N!$ is $\floor{\log N!}+1=\Theta(\log N!)$. We show
$\log N!=\Theta(N\log N)$ (Exercise 1.4).
\begin{itemize}[itemsep=0.2em]
\item Upper bound: $N!\le N^N$, so $\log N!\le N\log N$.
\item Lower bound: at least $N/2$ of the factors $1,\dots,N$ are $\ge N/2$, so $N!\ge (N/2)^{N/2}$ and
  $\log N!\ge\frac N2(\log N-1)\ge\frac N4\log N$ for $N\ge4$.
\end{itemize}
Since $N=\Theta(2^n)$ and $\log N=\Theta(n)$, the length of $N!$ is
\[
  \Theta(N\log N)=\boxed{\Theta(n\,2^n)}\text{ bits.}
\]

\textbf{(b)} Algorithm: multiply the numbers one by one.

\begin{algorithm}[H]
\caption{Computing $N!$}
\begin{algorithmic}[1]
\Function{Factorial}{$N$}
  \State $f\gets1$
  \For{$i\gets2$ \textbf{to} $N$}
    \State $f\gets f\cdot i$ \Comment{invariant: after this line, $f=i!$}
  \EndFor
  \State \Return $f$
\EndFunction
\end{algorithmic}
\end{algorithm}

\emph{Correctness} is immediate from the loop invariant $f=i!$.

\emph{Running time.} Multiplying a $b$-bit number by an $m$-bit number with $m\le b$ by the
grade-school method adds up $m$ shifted copies of the $b$-bit number, each addition costing $O(b+m)=O(b)$;
so it takes $O(bm)$ time.

In iteration $i$, the current value $f=(i-1)!\le N!$ has $b=O(\log N!)=O(N\log N)=O(Nn)$ bits by part (a),
and $i\le N$ has at most $n$ bits. Hence each iteration costs $O(Nn\cdot n)=O(Nn^2)$ (the loop counter
update is only $O(n)$). With $N-1$ iterations, the total time is
\[
  O(N\cdot Nn^2)=O(N^2n^2)=\boxed{O(n^2\,4^n)}.
\]
This bound is tight for this algorithm: in each of the last $N/2$ iterations
($N/2 < i \le N$), $f=(i-1)!\ge (N/2)!$ already has $\Omega(N\log N)=\Omega(Nn)$ bits, and $i$ has at least
$n-1$ bits, so these iterations alone cost $\Omega(N\cdot Nn\cdot n)$ under the grade-school method.

\begin{remark}
The running time is exponential in the input length $n$. This is unavoidable for \emph{any} algorithm:
by (a), merely writing down the output takes $\Theta(n2^n)$ time.
\end{remark}

%======================================================================
\section*{Exercise 1.35}

\textbf{(a)} $x$ is its own inverse modulo $p$ iff $x^2\equiv1\pmod p$, i.e.\ $p\mid x^2-1=(x-1)(x+1)$.
Since $p$ is prime, this holds iff $p\mid x-1$ or $p\mid x+1$, i.e.\ $x\equiv1$ or $x\equiv-1\pmod p$.
For $1\le x<p$, the self-inverse numbers are therefore exactly
\[
  x=1\quad\text{and}\quad x=p-1
\]
(for $p=2$ these coincide, leaving only $x=1$).

\textbf{(b)} If $p=2$: $(p-1)!=1\equiv-1\pmod 2$. Now let $p$ be an odd prime and
$S=\{2,3,\dots,p-2\}$.

Every $x\in S$ has a unique inverse $x^{-1}\in\{1,\dots,p-1\}$ modulo $p$. We claim $x^{-1}\in S$ and
$x^{-1}\neq x$:
\begin{itemize}[itemsep=0.2em]
\item If $x^{-1}=1$, then $x=1$; if $x^{-1}=p-1$, then $x=(p-1)^{-1}=p-1$ (as $p-1$ is self-inverse).
  Both contradict $x\in S$, so $x^{-1}\in S$.
\item $x^{-1}\neq x$ by part (a), since $x\notin\{1,p-1\}$.
\end{itemize}
Moreover $(x^{-1})^{-1}=x$. Hence $S$ splits into disjoint pairs $\{x,x^{-1}\}$, each with product
$\equiv1\pmod p$, and so $\prod_{x\in S}x\equiv1\pmod p$ (for $p=3$, $S=\emptyset$ and the empty product is $1$).
Therefore
\[
  (p-1)!=1\cdot\Bigl(\prod_{x\in S}x\Bigr)\cdot(p-1)\equiv1\cdot1\cdot(-1)=-1\pmod p .
\]

\textbf{(c)} Wilson's theorem concerns $N\ge2$ (modulo $1$ all integers are congruent, so $N=1$ is excluded).
Let $N\ge2$ be composite and write $N=ab$ with $1<a<N$. Since $a\le N-1$, $a$ is one of the factors of $(N-1)!$, so
$a\mid(N-1)!$. Also $a\mid N$. Hence
\[
  d=\gcd(N,(N-1)!)\ge a>1 .
\]
Suppose, for contradiction, that $(N-1)!\equiv-1\pmod N$, i.e.\ $(N-1)!=kN-1$ for some integer $k$. Then
$d\mid N$ and $d\mid(N-1)!$ imply $d\mid kN-(N-1)!=1$, so $d=1$, a contradiction.
Thus $(N-1)!\not\equiv-1\pmod N$.

\textbf{(d)} To use the test we must compute $(N-1)!\bmod N$. For an $n$-bit $N$ the straightforward way
performs $N-2=\Theta(2^n)$ modular multiplications (each $O(n^2)$), which is exponential in the input size $n$.
Unlike $a^{N-1}\bmod N$ in Fermat's test, which repeated squaring computes with only $O(n)$
multiplications because all factors are equal, the factors $1,2,\dots,N-1$ of $(N-1)!$ are all different,
and no polynomial-time algorithm for computing $(N-1)!\bmod N$ is known. So Wilson's theorem, although an exact
characterization of primes, does not yield an efficient primality test.

\end{document}
